\documentclass[10pt,letterpaper]{article}
\usepackage[utf8]{inputenc}
\usepackage[spanish,es-tabla,es-noquoting]{babel}
\usepackage{geometry}
\geometry{margin=0.64in, top=0.74in, bottom=0.65in, headheight=36pt, headsep=14pt}
\usepackage{amsmath,amssymb}
\usepackage{xcolor}
\usepackage{graphicx}
\usepackage{tabularx}
\usepackage{booktabs}
\usepackage{colortbl}
\usepackage{array}
\usepackage{enumitem}
\usepackage{fancyhdr}
\usepackage{titlesec}
\usepackage{lastpage}
\usepackage{setspace}
\usepackage{listings}
\usepackage{tcolorbox}

% Tipografía institucional Icesi (Sans-Serif Helvetica)
\usepackage{helvet}
\renewcommand{\familydefault}{\sfdefault}

% Interlineado aumentado un 30% (setstretch 1.38)
\setstretch{1.38}

% Altura de celdas en tablas
\renewcommand{\arraystretch}{1.25}

% Paleta de Colores Oficial de los Programas Icesi
\definecolor{icesiBlue}{HTML}{5454E9}       % Azul vibrante institucional
\definecolor{icesiOrange}{HTML}{F0681A}     % Naranja acento
\definecolor{icesiGreen}{HTML}{4CB978}      % Verde acento
\definecolor{icesiPurple}{HTML}{865CF0}     % Púrpura acento
\definecolor{icesiDark}{HTML}{1A1A1A}       % Texto principal oscuro
\definecolor{icesiMuted}{HTML}{6B7280}      % Gris metadatos
\definecolor{icesiBorder}{HTML}{D1D5DB}     % Líneas y bordes
\definecolor{icesiRowLine}{HTML}{E5E7EB}    % Líneas de tabla
\definecolor{icesiLight}{HTML}{F9FAFB}      % Fondo suave
\definecolor{codeBackground}{HTML}{F3F4F6}  % Fondo código

% Configuración de Listings
\lstdefinelanguage{JavaScript}{
  morekeywords={typeof, new, true, false, catch, function, return, null, switch, var, if, in, while, do, else, case, break, const, let, async, await, import, export, from, default, try, throw, class, extends, super},
  sensitive=true,
  morecomment=[l]{//},
  morecomment=[s]{/*}{*/},
  morestring=[b]',
  morestring=[b]"
}

\lstdefinestyle{jsstyle}{
    language=JavaScript,
    backgroundcolor=\color{codeBackground},
    commentstyle=\color{icesiMuted}\itshape,
    keywordstyle=\color{icesiBlue}\bfseries,
    stringstyle=\color{icesiOrange},
    basicstyle=\ttfamily\scriptsize,
    breaklines=true,
    frame=single,
    rulecolor=\color{icesiBorder},
    showstringspaces=false
}

\lstdefinestyle{bashstyle}{
    backgroundcolor=\color{codeBackground},
    commentstyle=\color{icesiMuted}\itshape,
    keywordstyle=\color{icesiBlue}\bfseries,
    basicstyle=\ttfamily\scriptsize,
    breaklines=true,
    frame=single,
    rulecolor=\color{icesiBorder},
    language=bash,
    showstringspaces=false
}

\lstset{
  inputencoding=utf8,
  extendedchars=true,
  literate=%
    {á}{{\'a}}1 {é}{{\'e}}1 {í}{{\'i}}1 {ó}{{\'o}}1 {ú}{{\'u}}1
    {Á}{{\'A}}1 {É}{{\'E}}1 {Í}{{\'I}}1 {Ó}{{\'O}}1 {Ú}{{\'U}}1
    {ñ}{{\~n}}1 {Ñ}{{\~N}}1
}

% =========================================================================
% ENCABEZADO Y PIE DE PÁGINA ESTILO PROGRAMA DE CURSO ICESI
% =========================================================================
\pagestyle{fancy}
\fancyhf{}
\fancyhead[L]{%
  \begin{minipage}[c]{2.8cm}
    \includegraphics[height=0.58cm]{logo-icesi.png}
  \end{minipage}%
  \hspace{0.15cm}%
  {\color{icesiBorder}\vrule width 0.8pt height 0.60cm}%
  \hspace{0.2cm}%
  \begin{minipage}[c]{10.5cm}
    \fontsize{7.5}{9.5}\selectfont\textbf{\color{icesiDark}UNIVERSIDAD ICESI} \\
    \fontsize{7}{9}\selectfont\textbf{\color{icesiDark}GUÍA DE TALLER PRÁCTICO: COMPUTACIÓN EN INTERNET I} \\
    \fontsize{6.5}{8.5}\selectfont\color{icesiMuted}Período: 2026-2 $\bullet$ NRC: 12378 / 12377 $\bullet$ Modalidad: Equipos de 2 Estudiantes
  \end{minipage}%
}
\fancyhead[R]{\fontsize{7.5}{9.5}\selectfont\color{icesiMuted}Pág. \thepage\ de \pageref{LastPage}}

% Barra horizontal multicolor
\renewcommand{\headrule}{%
  \vspace{0.08cm}%
  \noindent%
  \textcolor{icesiBlue}{\rule{0.60\textwidth}{1.8pt}}%
  \textcolor{icesiOrange}{\rule{0.18\textwidth}{1.8pt}}%
  \textcolor{icesiGreen}{\rule{0.11\textwidth}{1.8pt}}%
  \textcolor{icesiPurple}{\rule{0.11\textwidth}{1.8pt}}%
}
\renewcommand{\footrulewidth}{0.4pt}
\fancyfoot[L]{\fontsize{7.5}{9.5}\selectfont\color{icesiMuted}Computación en Internet I (09810 - TIC) $\bullet$ Taller: DOM, Componentes, Eventos y Asincronía en Vanilla JS}
\fancyfoot[R]{\fontsize{7.5}{9.5}\selectfont\color{icesiMuted}Universidad Icesi $\bullet$ DCSI}

% =========================================================================
% COMANDOS DE SECCIONES CON BADGE ESTILO PROGRAMA ICESI
% =========================================================================
\newcommand{\sectionBadge}[1]{%
  \refstepcounter{section}%
  \vspace{0.24cm}%
  \noindent%
  \colorbox{icesiBlue}{\color{white}\rule[-0.08cm]{0pt}{0.42cm}\hspace{0.16cm}\textbf{\normalsize \thesection}\hspace{0.16cm}}%
  \hspace{0.18cm}%
  {\fontsize{9.5}{12}\selectfont\textbf{\color{icesiDark}#1}}\\[-0.10cm]%
  \noindent{\color{icesiDark}\rule{\textwidth}{0.8pt}}\vspace{0.10cm}%
}

\titleformat{\subsection}
  {\color{icesiBlue}\normalfont\normalsize\bfseries}
  {\thesubsection}{0.4em}{}
\titlespacing*{\subsection}{0pt}{0.6ex plus 0.2ex minus .1ex}{0.2ex plus .1ex}

\newcolumntype{C}[1]{>{\centering\arraybackslash}p{#1}}
\newcolumntype{L}[1]{>{\raggedright\arraybackslash}p{#1}}
\newcolumntype{Y}{>{\centering\arraybackslash}X}

\begin{document}
\shorthandoff{"}

% =========================================================================
% TÍTULO Y PRESENTACIÓN INICIAL
% =========================================================================
\noindent
{\fontsize{8.5}{11}\selectfont\textbf{\color{icesiBlue}TALLER PRÁCTICO FORMATIVO: ARQUITECTURA WEB CLIENTE · SESIONES 26 A 28}}\\[0.06cm]
{\fontsize{13.5}{17}\selectfont\textbf{Manipulación del DOM, Componentes Reutilizables, Eventos, Promesas y Animaciones en Vanilla JS}}\\[0.10cm]
{\color{icesiRowLine}\hrule height 0.6pt}\vspace{0.10cm}

\noindent
{\color{icesiBlue}\textbf{\rule{2.5pt}{8.5pt}}}~\textbf{\fontsize{8.5}{11}\selectfont FACULTAD DE INGENIERÍA, DISEÑO Y CIENCIAS APLICADAS} \hfill
{\color{icesiOrange}\textbf{\rule{2.5pt}{8.5pt}}}~\textbf{\fontsize{8.5}{11}\selectfont DEPARTAMENTO DE COMPUTACIÓN Y SISTEMAS INTELIGENTES}\\[0.06cm]
{\color{icesiRowLine}\hrule height 0.6pt}\vspace{0.10cm}

% =========================================================================
% TABLA DE IDENTIFICACIÓN
% =========================================================================
\noindent
{\color{icesiBlue}\textbf{\rule{2.5pt}{8.5pt}}}~\textbf{\fontsize{8.5}{11}\selectfont INFORMACIÓN GENERAL} \hfill {\fontsize{7.5}{9}\selectfont\color{icesiMuted}\textbf{IDENTIFICACIÓN DEL CURSO}}\\[-0.22cm]
{\color{icesiBorder}\hrule height 0.5pt}\vspace{0.05cm}

\noindent
\small
\begin{tabularx}{\textwidth}{>{\bfseries\color{icesiDark}\fontsize{7.5}{9.5}\selectfont}L{2.6cm}X!{\color{icesiBlue}\vrule width 0.8pt}>{\bfseries\color{icesiDark}\fontsize{7.5}{9.5}\selectfont}L{2.6cm}X}
CÓDIGO: & 09810 - TIC & PERÍODO: & 202620 (Semestre 2026-2) \\ \hline
\rowcolor{icesiLight}
NRC / GRUPO: & 12378 / 12377 (Grupos 001 y 003) & CARÁCTER: & Práctico Formativo y Aplicado \\ \hline
MODALIDAD: & Equipos de 2 Estudiantes & DURACIÓN: & 2 Horas en Clase + Cierre Autónomo \\ \hline
\rowcolor{icesiLight}
POLÍTICA IAG: & Nivel 2: Asistencia de Consulta y Depuración & PONDERACIÓN: & Formativo (Alineación RAA 3 / Preparación Parcial 3) \\ \hline
ESTUDIANTE 1: & \rule{4.5cm}{0.4pt} & CÓDIGO: & \rule{2.2cm}{0.4pt} \\ \hline
\rowcolor{icesiLight}
ESTUDIANTE 2: & \rule{4.5cm}{0.4pt} & CÓDIGO: & \rule{2.2cm}{0.4pt} \\ \hline
\end{tabularx}

\vspace{0.10cm}

% =========================================================================
% SECCION 1: CONTEXTO Y STRICT MODE
% =========================================================================
\sectionBadge{Contexto: Modo Estricto y Manipulación Reactiva del DOM}

En el desarrollo web moderno, las interfaces de usuario no se conciben como documentos estáticos impresos, sino como \textbf{aplicaciones reactivas en memoria}. Comprender las bases de Vanilla JavaScript permite al ingeniero construir interfaces de alto rendimiento, dominar el ciclo de vida del navegador y comprender con precisión qué ocurre por debajo de cualquier framework contemporáneo.

\begin{tcolorbox}[colback=icesiOrange!10,colframe=icesiOrange,arc=3pt,title=\textbf{Directiva de Seguridad: Modo Estricto ("use strict")}]
\fontsize{7.5}{9.5}\selectfont
El Modo Estricto transforma errores silenciosos y conductas inseguras en excepciones explícitas inmediatas (\texttt{throw error}). Prohíbe variables globales accidentales (\texttt{x = 10} sin declarar lanza \texttt{ReferenceError}), asigna \texttt{undefined} a \texttt{this} en funciones sueltas y previene colisiones en compiladores JIT (V8). En módulos ES6 nativos (\texttt{<script type="module">}), el modo estricto está habilitado de forma automática e irrevocable.
\end{tcolorbox}

\subsection{Mutación Limpia vs. Antipatrones Inseguros}
El uso indiscriminado de \texttt{element.innerHTML = "..."} destruye listeners previos, genera vulnerabilidades de inyección de código (\textbf{Cross-Site Scripting - XSS}) y satura el motor gráfico con múltiples ciclos de \emph{Reflow} (cálculo geométrico) y \emph{Repaint} (redibujado de píxeles).

\begin{lstlisting}[style=jsstyle]
"use strict";

// ANTIPATRON: Inseguro contra ataques XSS y costoso en render
function renderItemInseguro(tarea, contenedor) {
  contenedor.innerHTML += `<div class="item">${tarea.titulo}</div>`;
}

// PATRON RECOMENDADO: Nodos en memoria con sanitizacion nativa
function renderItemSeguro(tarea, contenedor) {
  const item = document.createElement("article");
  item.className = "task-item";
  item.dataset.id = tarea.id;

  const title = document.createElement("span");
  title.textContent = tarea.titulo; // 100% inmune a XSS

  const btn = document.createElement("button");
  btn.className = "btn-delete";
  btn.textContent = "Eliminar";

  item.append(title, btn);
  contenedor.appendChild(item);
}
\end{lstlisting}

% =========================================================================
% SECCION 2: ARQUITECTURA DE COMPONENTES REUTILIZABLES
% =========================================================================
\sectionBadge{Arquitectura de Componentes Reutilizables en Vanilla JS}

Un componente funcional en Vanilla JavaScript es una función pura que recibe un objeto de propiedades de entrada (\textbf{props}) y retorna un nodo del DOM completamente configurado.

\begin{itemize}[leftmargin=*,itemsep=2pt,topsep=2pt]
  \item \textbf{Paso de Props Descendente (\textit{Props Down}):} Los componentes padres transmiten datos hacia los componentes hijos mediante argumentos.
  \item \textbf{Notificación Ascendente (\textit{Callbacks Up}):} Los componentes hijos nunca mutan el estado central directamente; notifican al padre invocando funciones callback pasadas en sus props (ej. \texttt{onDelete(id)}, \texttt{onToggle(id)}).
  \item \textbf{Estado Centralizado en el Componente Raíz (\texttt{App}):} Toda la lógica de negocio reside en el componente raíz, garantizando una única fuente de verdad (\textit{Single Source of Truth}).
\end{itemize}

% =========================================================================
% SECCION 3: EVENTOS, PROPAGACION Y DELEGACION
% =========================================================================
\sectionBadge{Propagación de Eventos y el Patrón de Delegación}

Cuando ocurre una interacción física en el navegador, el evento recorre tres fases formales:
\begin{enumerate}[label=\arabic*.]
  \item \textbf{Fase de Captura (\textit{Capturing}):} El evento desciende desde \texttt{window} y \texttt{document} hasta alcanzar el nodo objetivo.
  \item \textbf{Fase de Destino (\textit{Target}):} Se alcanza el nodo sobre el que se originó el clic (\texttt{event.target}).
  \item \textbf{Fase de Burbujeo (\textit{Bubbling}):} El evento asciende hacia los ancestros por defecto (\texttt{bubbles: true}).
\end{enumerate}

\begin{tcolorbox}[colback=icesiBlue!10,colframe=icesiBlue,arc=3pt,title=\textbf{Eficiencia de la Delegación de Eventos}]
\fontsize{7.5}{9.5}\selectfont
En lugar de asociar 500 listeners a 500 botones individuales en memoria RAM, registramos \textbf{un único listener en el contenedor padre} (\texttt{<ul>}) y filtramos con \texttt{event.target.closest(".btn-delete")}. Esto soporta automáticamente nuevos elementos creados dinámicamente sin fugas de memoria.
\end{tcolorbox}

% =========================================================================
% SECCION 4: ASINCRONIA, PROMESAS Y EL EVENT LOOP
% =========================================================================
\sectionBadge{Asincronía: Promesas, Async/Await y Prioridad de Microtareas}

El motor JavaScript ejecuta en un único hilo. Las operaciones asíncronas no bloqueantes se gestionan mediante el \textbf{Event Loop}:
\begin{itemize}[leftmargin=*,itemsep=2pt]
  \item \textbf{Microtasks (Alta Prioridad):} Las Promesas resueltas (\texttt{.then()}, \texttt{async/await}) entran a la cola de microtareas y se ejecutan inmediatamente apenas se vacía el Call Stack.
  \item \textbf{Macrotasks (Baja Prioridad):} Los temporizadores (\texttt{setTimeout}) esperan a que la cola de microtareas esté totalmente vacía antes de ejecutarse.
\end{itemize}

% =========================================================================
% SECCION 5: ANIMACIONES WEB Y GRAFICOS VECTORIALES SVG
% =========================================================================
\sectionBadge{Animaciones con Web Animations API y Trazado SVG}

Para dinamizar la interfaz sin degradar el rendimiento, utilizamos:
\begin{itemize}[leftmargin=*,itemsep=2pt]
  \item \textbf{Web Animations API (WAAPI):} \texttt{elem.animate([...], \{duration: 300\})} ejecutado sobre \texttt{transform} y \texttt{opacity} en GPU.
  \item \textbf{Animación Vectorial SVG:} Manipulación de \texttt{stroke-dashoffset} y atributos geométricos interactivos.
\end{itemize}

% =========================================================================
% SECCION 6: RETO PRACTICO Y RUBRICA
% =========================================================================
\sectionBadge{Reto Práctico: Dashboard DevPulse y Rúbrica Oficial}

Cada equipo construirá la aplicación modular de observabilidad \textbf{DevPulse} en Vanilla JavaScript puro (ES6+) partiendo de un contenedor raíz limpio (\texttt{<div id='root'></div>}).

\textbf{Requerimientos Principales del Proyecto:}
\begin{itemize}[leftmargin=*,itemsep=1.5pt,topsep=2pt]
  \item \textbf{Renderizado Dinámico y Seguro:} Renderizar catálogo de microservicios (\texttt{Auth}, \texttt{Payments}, \texttt{DB}) mediante \texttt{createElement()} y \texttt{textContent}, 100\% inmune a XSS. Prohibido \texttt{innerHTML}.
  \item \textbf{Arquitectura de Componentes Funcionales:} Desacoplamiento en componentes puros (\texttt{ServiceCard}, \texttt{MetricGaugeSVG}, \texttt{FilterBar}) con \textit{Props Down} y \textit{Callbacks Up}.
  \item \textbf{Delegación de Eventos:} Un único listener delegado sobre el contenedor \texttt{\#services-grid} gestionando acciones de reinicio, cortacircuitos y eliminación con \texttt{closest('[data-action]')}.
  \item \textbf{Asincronía y Simulación de Fallos:} Servicio de sondeo con Promesas y \texttt{async/await}, controlando estados \textit{Pending} (carga visual), \textit{Fulfilled} y \textit{Rejected} (captura de errores 500/503 con \texttt{try...catch...finally}).
  \item \textbf{Métricas SVG y Web Animations API:} Gráficos radiales vectoriales SVG con actualización de \texttt{stroke-dashoffset} y animaciones a 60 FPS con \texttt{element.animate()} en GPU.
\end{itemize}

\vspace{0.10cm}

\noindent
\textbf{Rúbrica de Evaluación Modularizada (Escala 0.0 a 5.0):}
\begin{center}
\small
\begin{tabularx}{\textwidth}{L{3.4cm}C{1.3cm}X}
\toprule
\textbf{Criterio} & \textbf{Peso} & \textbf{Indicador de Logro} \\
\midrule
\textbf{Manipulación del DOM y Erradicación XSS} & 20\% & Creación 100\% programática con \texttt{createElement()}, inserción segura con \texttt{textContent} y reemplazo atómico con \texttt{replaceChildren()}. Cero XSS. \\
\rowcolor{icesiLight}
\textbf{Arquitectura de Componentes y Estado} & 20\% & Componentes funcionales puros desacoplados, flujo Props Down / Callbacks Up y estado centralizado único sin variables globales. \\
\textbf{Delegación de Eventos y Memoria} & 20\% & Un único listener padre delegado con filtrado mediante \texttt{closest()} y \texttt{matches()}, soportando mutaciones dinámicas sin fugas. \\
\rowcolor{icesiLight}
\textbf{Asincronía con Promesas y Async/Await} & 20\% & Control riguroso de estados (Pending, Fulfilled, Rejected), manejo de excepciones con \texttt{try...catch...finally} y simulación de latencia de red. \\
\textbf{Animaciones WAAPI, SVG y Modo Estricto} & 20\% & Transiciones fluidas a 60 FPS con Web Animations API, medidor radial interactivo SVG y cumplimiento estricto de \texttt{"use strict";}. \\
\bottomrule
\end{tabularx}
\end{center}

\vspace{0.10cm}
\begin{tcolorbox}[colback=icesiGreen!10,colframe=icesiGreen,arc=3pt,title=\textbf{Portal Interactivo de Aprendizaje Web}]
\fontsize{7.5}{9.5}\selectfont
Este taller cuenta con un simulador interactivo completo en el portal Docusaurus del curso (\texttt{portal/docs/talleres/taller-dom-vanillajs-componentes}), con un Sandbox de código en vivo y simuladores visuales en tiempo real del Event Loop, de la propagación de eventos y del ciclo de vida de componentes.
\end{tcolorbox}

\end{document}
